% GENERATED FILE. Edit canonical lessons, then run npm run generate:books.
% canonical-source-hash: fnv1a64:4def3607834ac63b
% canonical-lessons: MR-C217-har, MR-C217-kanatle, MR-C217-bajarpeth, MR-C217-bandar, MR-C217-vihir

\chapter{Necklace, Earrings, Market quarter, Harbour, Well}
\label{ch:mr-necklace-earrings-market-quarter-harbour-well}
\hlchaptermodality{217}

\begin{chapteropening}
\textbf{By the end of this chapter:} \emph{I can say a necklace, earrings, a market quarter, a harbour and a well.}

Complete the last lesson of chapter 217: I can say a necklace, earrings, a market quarter, a harbour and a well.
\end{chapteropening}

\section[\texorpdfstring{\mr{हार} (hār)}{हार (hār)}]{\mr{हार} (hār) — a necklace, a garland}

\label{lesson:MR-C217-har}



\begin{quote}
Before the new one: say the Marathi for turmeric, then the Marathi for cumin.
\end{quote}

\subsection*{You'll want to know: \mr{हार}}
\textbf{\mr{हार}} — \emph{hār} — \textquotedblleft{}a necklace, a garland\textquotedblright{}.

\begin{cousinweb}
\textbf{\mr{हार}} (\emph{hār}) is a tadbhava, a word inherited from Sanskrit and reshaped by sound change. It comes from \emph{hāra}, \textquotedblleft{}necklace\textquotedblright{} (CDIAL 14059). CDIAL glosses the Marathi word \textquotedblleft{}garland; line\textquotedblright{}.
\end{cousinweb}

\begin{grammarlens}[title={What you've built}]
One more word for this chapter.
\end{grammarlens}

\subsection*{Your turn}
\begin{itemize}
\raggedright
  \item \emph{Say it:} \emph{hār}
  \item \emph{Say it:} \emph{hār}, once more
  \item \emph{Say it:} say \emph{hār}
  \item \emph{Recall:} say the Marathi for turmeric, then the Marathi for cumin, then say \emph{hār} again
\end{itemize}

\begin{tcolorbox}[breakable,colback=teal!4,colframe=teal!35!black,title={Before you move on}]
What does \mr{हार} mean? (\textquotedblleft{}A necklace, a garland\textquotedblright{}.) Say it once more.
\end{tcolorbox}
\section[\texorpdfstring{\mr{कानातले} (kānātle)}{कानातले (kānātle)}]{\mr{कानातले} (kānātle) — earrings}

\label{lesson:MR-C217-kanatle}



\begin{quote}
Before the new one: say the Marathi for cumin, then the Marathi for a necklace.
\end{quote}

\subsection*{You'll want to know: \mr{कानातले}}
\textbf{\mr{कानातले}} — \emph{kānātle} — \textquotedblleft{}earrings\textquotedblright{}.

\begin{cousinweb}
\textbf{\mr{कानातले}} (\emph{kānātle}) is \textquotedblleft{}the in-ear ones\textquotedblright{}: \textbf{\mr{कान}}, \textquotedblleft{}ear\textquotedblright{}, from Sanskrit \emph{karṇa}, \textquotedblleft{}ear\textquotedblright{} (CDIAL 2830), with \textbf{-\mr{त}}, \textquotedblleft{}in\textquotedblright{}, and the ending \emph{-le}.
\end{cousinweb}

\begin{grammarlens}[title={What you've built}]
One more word for this chapter.
\end{grammarlens}

\subsection*{Your turn}
\begin{itemize}
\raggedright
  \item \emph{Say it:} \emph{kānātle}
  \item \emph{Say it:} \emph{kānātle}, once more
  \item \emph{Say it:} say \emph{kānātle}
  \item \emph{Recall:} say the Marathi for cumin, then the Marathi for a necklace, then say \emph{kānātle} again
\end{itemize}

\begin{tcolorbox}[breakable,colback=teal!4,colframe=teal!35!black,title={Before you move on}]
What does \mr{कानातले} mean? (\textquotedblleft{}Earrings\textquotedblright{}.) Say it once more.
\end{tcolorbox}
\section[\texorpdfstring{\mr{बाजारपेठ} (bājārpeṭh)}{बाजारपेठ (bājārpeṭh)}]{\mr{बाजारपेठ} (bājārpeṭh) — a market quarter}

\label{lesson:MR-C217-bajarpeth}



\begin{quote}
Before the new one: say the Marathi for a necklace, then the Marathi for earrings.
\end{quote}

\subsection*{You'll want to know: \mr{बाजारपेठ}}
\textbf{\mr{बाजारपेठ}} — \emph{bājārpeṭh} — \textquotedblleft{}a market quarter\textquotedblright{}.

\begin{cousinweb}
\textbf{\mr{बाजारपेठ}} (\emph{bājārpeṭh}) joins two words for market: \emph{bājār}, usually traced to Persian \emph{bāzār}, and \emph{peṭh}, \textquotedblleft{}town quarter, market\textquotedblright{}, which CDIAL lists under Sanskrit \emph{pratiṣṭhā}, \textquotedblleft{}foundation; homestead\textquotedblright{} (CDIAL 8608).
\end{cousinweb}

\begin{grammarlens}[title={What you've built}]
One more word for this chapter.
\end{grammarlens}

\subsection*{Your turn}
\begin{itemize}
\raggedright
  \item \emph{Say it:} \emph{bājārpeṭh}
  \item \emph{Say it:} \emph{bājārpeṭh}, once more
  \item \emph{Say it:} say \emph{bājārpeṭh}
  \item \emph{Recall:} say the Marathi for a necklace, then the Marathi for earrings, then say \emph{bājārpeṭh} again
\end{itemize}

\begin{tcolorbox}[breakable,colback=teal!4,colframe=teal!35!black,title={Before you move on}]
What does \mr{बाजारपेठ} mean? (\textquotedblleft{}A market quarter\textquotedblright{}.) Say it once more.
\end{tcolorbox}
\section[\texorpdfstring{\mr{बंदर} (bandar)}{बंदर (bandar)}]{\mr{बंदर} (bandar) — a harbour}

\label{lesson:MR-C217-bandar}



\begin{quote}
Before the new one: say the Marathi for earrings, then the Marathi for a market quarter.
\end{quote}

\subsection*{You'll want to know: \mr{बंदर}}
\textbf{\mr{बंदर}} — \emph{bandar} — \textquotedblleft{}a harbour\textquotedblright{}.

\begin{cousinweb}
\textbf{\mr{बंदर}} (\emph{bandar}), \textquotedblleft{}harbour\textquotedblright{}, is a loan, usually traced to Persian \emph{bandar}, \textquotedblleft{}port\textquotedblright{}.
\end{cousinweb}

\begin{grammarlens}[title={What you've built}]
One more word for this chapter.
\end{grammarlens}

\subsection*{Your turn}
\begin{itemize}
\raggedright
  \item \emph{Say it:} \emph{bandar}
  \item \emph{Say it:} \emph{bandar}, once more
  \item \emph{Say it:} say \emph{bandar}
  \item \emph{Recall:} say the Marathi for earrings, then the Marathi for a market quarter, then say \emph{bandar} again
\end{itemize}

\begin{tcolorbox}[breakable,colback=teal!4,colframe=teal!35!black,title={Before you move on}]
What does \mr{बंदर} mean? (\textquotedblleft{}A harbour\textquotedblright{}.) Say it once more.
\end{tcolorbox}
\section[\texorpdfstring{\mr{विहीर} (vihīr)}{विहीर (vihīr)}]{\mr{विहीर} (vihīr) — a well}

\label{lesson:MR-C217-vihir}



\begin{quote}
Before the new one: say the Marathi for a market quarter, then the Marathi for a harbour.
\end{quote}

\subsection*{You'll want to know: \mr{विहीर}}
\textbf{\mr{विहीर}} — \emph{vihīr} — \textquotedblleft{}a well\textquotedblright{}.

\begin{cousinweb}
The origin of \textbf{\mr{विहीर}} (\emph{vihīr}) is not settled. Neither CDIAL nor DEDR, searched by form and by meaning, has an entry for it.
\end{cousinweb}

\begin{grammarlens}[title={What you've built}]
One more word for this chapter.
\end{grammarlens}

\subsection*{Your turn}
\begin{itemize}
\raggedright
  \item \emph{Say it:} \emph{vihīr}
  \item \emph{Say it:} \emph{vihīr}, once more
  \item \emph{Say it:} say \emph{vihīr}
  \item \emph{Recall:} say the Marathi for a market quarter, then the Marathi for a harbour, then say \emph{vihīr} again
\end{itemize}

\begin{tcolorbox}[breakable,colback=teal!4,colframe=teal!35!black,title={Before you move on}]
What does \mr{विहीर} mean? (\textquotedblleft{}A well\textquotedblright{}.) Say it once more.
\end{tcolorbox}
